% Generated draft: requires mathematical and source review.
\documentclass[11pt]{article}
\usepackage[margin=25mm]{geometry}
\usepackage{fontspec}
\setmainfont{Latin Modern Roman}
\usepackage{amsmath,amssymb,mathtools,unicode-math}
\setmathfont{Latin Modern Math}
\usepackage{enumitem,needspace,xurl,hyperref}
\setlength{\parindent}{0pt}
\setlength{\parskip}{5pt}
\setlength{\emergencystretch}{3em}
\newcommand{\N}{\mathbb{N}}
\newcommand{\Z}{\mathbb{Z}}
\newcommand{\Q}{\mathbb{Q}}
\newcommand{\R}{\mathbb{R}}
\newcommand{\C}{\mathbb{C}}
\newcommand{\F}{\mathbb{F}}
\newcommand{\A}{\mathbb{A}}
\newcommand{\PP}{\mathbb{P}}
\newcommand{\E}{\mathbb{E}}
\newcommand{\eps}{\varepsilon}
\newcommand{\dd}{\,\mathrm{d}}
\begin{document}
\section*{UnsolvedMath ID: 20003140}
\textbf{Source problem code:} catalogue entry
\subsection{Definitions and mathematical statement}
The source asks whether every finite-dimensional Alexandrov space is a limit of Riemannian manifolds with \(\sec\ge K\). The role of \(K\) and the dimensions must be fixed: preserving the space's prescribed synthetic lower bound, allowing a weaker common bound, requiring equal dimension, allowing uniformly bounded higher dimensions, and allowing dimensions tending to infinity are different statements. For a compact target use Gromov--Hausdorff convergence of compact Riemannian manifolds; the source does not specify boundary or pointed/noncompact variants.

The supplied aid considers
\[
X=\operatorname{Susp}_s(\R\PP^4),
\]
where \(\R\PP^4\) has its round curvature-one metric and the suspension is spherical. It reports that \(X\) is a compact five-dimensional Alexandrov space with optimal synthetic curvature lower bound one, but for every fixed \(\kappa>1/4\) there is no approximating sequence satisfying
\[
M_i\xrightarrow{\rm GH}X,\qquad \sec(M_i)\ge\kappa,\qquad \sup_i\dim M_i<\infty.
\]
The endpoint \(\kappa=1/4\) is not excluded by that reported strict threshold. In particular the same \(K=1\) cannot be retained in any bounded dimension.

The aid distinguishes two other conclusions. First, the Petersen--Wilhelm--Zhu quotient theorem reportedly supplies a fixed higher-dimensional approximation to this \(X\) with some weaker uniform lower bound; no numerical value or full theorem hypotheses are supplied. Thus this example is not an obstruction to every weaker-bound approximation in higher dimension. Second, local homology and Perelman stability rule out any same-dimensional approximation with any common lower bound. At a suspension pole \(v\),
\[
H_i(X,X\setminus\{v\};\Z)\cong\widetilde H_{i-1}(\R\PP^4;\Z),
\]
which is not the local homology of a five-manifold. The noncollapsed five-dimensional stability argument therefore cannot yield such a singular target from five-manifolds with a common sectional-curvature lower bound. The applicable compact/boundaryless stability hypotheses must be retained.

To isolate the quantified obstruction, define
\[
\Lambda_{\rm bd}(X)=\sup\{\kappa\in\R:\ \exists D<\infty,\ \exists M_i\xrightarrow{\rm GH}X,
\ \dim M_i\le D,\ \sec(M_i)\ge\kappa\}.
\]
The report gives \(\Lambda_{\rm bd}(X)\le1/4\), whereas the optimal Alexandrov bound is one. It does not compute \(\Lambda_{\rm bd}\) exactly. Any hypothetical same-bound \(K=1\) approximation without a dimension restriction would have to satisfy \(\dim M_i\to\infty\), since a bounded-dimensional subsequence would contradict the reported obstruction. Neither existence nor impossibility of that unrestricted sequence is established in the supplied summary. Accordingly the example refutes prescribed-bound bounded-dimensional formulations, rather than automatically settling every reading of the one-line question.
\subsection{Short English statement}
Ask whether Alexandrov spaces admit curvature-preserving Riemannian approximations, keeping dimension and the meaning of the bound explicit. The supplied suspension example obstructs same-bound bounded-dimensional smoothing, permits a weaker-bound higher-dimensional approximation, and leaves unbounded-dimensional same-bound approximation unestablished.
\subsection{Sources}
\begin{itemize}
\item Problem 15, AIM workshop \emph{Manifolds with nonnegative sectional curvature}, \url{https://aimath.org/WWN/nnsectcurvature/nnsectcurvature.pdf}; verification says the wording has no OCR ambiguity.
\item The research summary cites Petersen--Wilhelm--Zhu's quotient theorem and Perelman stability, without full bibliographic entries or a separately identified reference for the strict \(1/4\) bounded-dimension obstruction.
\item Local UnsolvedMath registry, permanent ID 20003140, exact code \texttt{catalogue entry}; canonical location \texttt{aim-topology-notes.json}, notes[227], dated 14 August 2026; statement labeled exact. Registry status solved and classification ``solved by you'' are supplied metadata of the prior research record, not a claim of a new proof in this draft.
\item Research attempt 1: counterexample, problem partially solved at the run, rigor label proved, novelty confidence low. The curvature-threshold and quotient account is a machine-generated research aid, not peer reviewed; independent verification is required before citation.
\end{itemize}
\end{document}
